% Relatório autossuficiente: não requer figuras ou arquivos de dados externos.
% Compilar duas vezes com pdfLaTeX para atualizar as referências internas.
\documentclass[11pt,a4paper]{article}
\usepackage[utf8]{inputenc}
\usepackage[T1]{fontenc}
\usepackage[brazil]{babel}
\usepackage[margin=2.3cm]{geometry}
\usepackage{booktabs,tabularx,array}
\usepackage{xurl}
\usepackage[hidelinks]{hyperref}
\urlstyle{tt}
\setlength{\parindent}{0pt}
\setlength{\parskip}{0.55em}
\setlength{\emergencystretch}{2em}
\renewcommand{\arraystretch}{1.15}
\newcommand{\hex}[1]{\texttt{0x#1}}

\title{Relatório de simulação pós-síntese com anotação SDF\\
RSA de 32 bits com interface AMBA APB\\
\large TSMC28 --- clock de 100 MHz}
\author{Projeto RSA\_APB\_32-bit}
\date{Execução de 5 de novembro de 2026, às 14h37min30s (UTC$-$3)}

\begin{document}
\maketitle

\begin{abstract}
A última simulação pós-síntese disponível do periférico RSA/APB concluiu
as 398 verificações previstas, com 398 acertos e nenhuma falha, utilizando
clock de 100 MHz e a netlist TSMC28 sintetizada às 14h36min09s.
O tempo simulado foi de 134,98~$\mu$s. A anotação alcançou 100\% dos atrasos
de caminho e 41,80\% dos checks temporais existentes nos modelos de células.
Todos os checks \texttt{SETUPHOLD} e \texttt{RECREM} fornecidos pelo SDF
foram anotados. Permaneceram avisos relacionados a checks temporais negativos
e à desativação de checks de largura de pulso. O resultado confirma o
funcionamento para os estímulos executados, com as limitações de timing
detalhadas neste documento.
\end{abstract}

\section{Identificação e condições da execução}

Este relatório foi elaborado a partir dos arquivos da execução identificada
abaixo, sem realizar nova síntese ou simulação. Os caminhos de referência
são relativos à raiz do repositório.

\textbf{Identificador da simulação:}\\
\path{amba_3_apb_rsa_apb_top_2026-11-05_14-37-30_767454205_tsmc28}

\textbf{Identificador da síntese utilizada:}\\
\path{amba_3_apb_rsa_apb_top_2026-11-05_14-36-09_068054729_tsmc28}

\begin{center}
\begin{tabularx}{\linewidth}{@{}lX@{}}
\toprule
Item & Configuração registrada \\
\midrule
Topo & \path{amba_3_apb_rsa_apb_top} \\
Instância anotada & \path{testbench.dut} \\
Simulador & Cadence Xcelium/Xrun 24.03-s013 \\
Síntese & Cadence Genus 22.18-s061\_1 \\
Biblioteca Liberty & \path{/tsmc28/tcbn28hpcplusbwp30p140tt1v25c.lib} \\
Modelos Verilog & \path{/tsmc28/tcbn28hpcplusbwp30p140.v} \\
Corner & TT, 1,0 V, 25~$^{\circ}$C \\
Clock & 10,0 ns / 100 MHz \\
Escalas de tempo & Testbench: 1 ns / 1 ps; SDF: 1 ps \\
Seleção SDF & \texttt{MAXIMUM} \\
Defines & \texttt{POST\_SYNTHESIS}, \texttt{NTC}, \texttt{RECREM} \\
Carga dos operandos & \texttt{LOAD\_GAP=1}, confirmado no log de elaboração do Genus \\
\bottomrule
\end{tabularx}
\end{center}

A síntese é lógica, sem LEF e sem parasitas extraídos de layout. O relatório
de área identifica \texttt{ZeroWireload}, em modo \texttt{segmented}.
A opção \texttt{MAXIMUM} seleciona o campo máximo do SDF; ela não altera
o corner de caracterização para um corner de pior caso. O cabeçalho do
SDF registra explicitamente 1,0 V e 25~$^{\circ}$C.

\section{Método de verificação}

O testbench \path{testbenchs/rsa_apb_top/testbench.sv} aplica transações APB,
compara os dados lidos com os valores esperados e contabiliza os resultados
em um scoreboard. Os estímulos APB são aplicados na borda de descida do
clock, afastados da borda de amostragem do DUT. Para a simulação pós-síntese,
a configuração do hardware está incorporada à netlist.

Foram executadas oito verificações iniciais, 30 rodadas de banco de
registradores com oito verificações por rodada e dez rodadas de RSA com
15 verificações por rodada:
\[
N_{\mathrm{checks}} = 8 + 30\times8 + 10\times15 = 398.
\]
Os testes RSA verificam carga e preservação dos operandos, conclusão da
operação, resultado da exponenciação modular, estado e interrupção.
O valor esperado é calculado pela função de referência de exponenciação
modular do testbench. Os limites configurados são 20 ciclos para uma
transação APB e 10\,000 consultas de estado por operação RSA.

\section{Resultados funcionais}

\begin{center}
\begin{tabular}{@{}lr@{}}
\toprule
Métrica & Resultado \\
\midrule
Verificações previstas / realizadas & 398 / 398 \\
Acertos / falhas & 398 / 0 \\
Rodadas do banco de registradores & 30 \\
Rodadas RSA & 10 \\
Tempo simulado & 134\,980 ns (134,98 $\mu$s) \\
Tempo de execução informado pelo Xrun & 3 s \\
Código de saída do Xrun / fluxo & 0 / 0 \\
Mensagens de erro ou fatal (\texttt{*E}, \texttt{*F}) & 0 \\
Ocorrências de \texttt{Timing violation} no log & 0 \\
\bottomrule
\end{tabular}
\end{center}

O log encerra com \texttt{TEST PASS: 398 checks, 0 errors} e término normal
por \texttt{\$finish}. A Tabela~\ref{tab:vectors} apresenta os dez vetores
executados. Em cada linha, o valor lido em \texttt{RESULT} coincidiu com o
valor esperado. Todos os operandos e resultados da tabela estão em
hexadecimal, com exceção do número da rodada e do expoente.

\begin{table}[htbp]
\centering
\small
\begin{tabular}{@{}rllrl l@{}}
\toprule
Rodada & Mensagem & Módulo & Expoente & Esperado = lido & Estado \\
\midrule
1  & \hex{00000041} & \hex{00000ca1} & 17 & \hex{00000ae6} & PASS \\
2  & \hex{0000007b} & \hex{00000ca1} & 7  & \hex{00000b34} & PASS \\
3  & \hex{00000000} & \hex{00000061} & 5  & \hex{00000000} & PASS \\
4  & \hex{00000001} & \hex{00000101} & 13 & \hex{00000001} & PASS \\
5  & \hex{0000002a} & \hex{00000061} & 0  & \hex{00000001} & PASS \\
6  & \hex{12345678} & \hex{0000fff1} & 1  & \hex{00006793} & PASS \\
7  & \hex{80000001} & \hex{fffffffb} & 2  & \hex{4000000b} & PASS \\
8  & \hex{ffffffff} & \hex{fffffffb} & 3  & \hex{00000040} & PASS \\
9  & \hex{075bcd15} & \hex{00010001} & 7  & \hex{00008a8e} & PASS \\
10 & \hex{fedcba98} & \hex{8000000b} & 17 & \hex{14adadfc} & PASS \\
\bottomrule
\end{tabular}
\caption{Vetores RSA e resultados observados no log da simulação.}
\label{tab:vectors}
\end{table}

Na última rodada, o resultado \hex{14adadfc} foi lido em 134\,830 ns.
As verificações finais confirmaram a preservação de mensagem, expoente
e módulo, o retorno de \texttt{STATUS} a zero e a desativação de IRQ
em 134\,980 ns. A taxa de acerto de 100\% refere-se às comparações
executadas; ela não representa uma medida de cobertura funcional ou de código.

\section{Anotação SDF e avisos}

\begin{center}
\begin{tabular}{@{}lrrr@{}}
\toprule
Categoria & Total no modelo & Anotados & Percentual \\
\midrule
Atrasos de caminho & 14\,871 & 14\,871 & 100,00\% \\
Checks temporais (total) & 17\,826 & 7\,452 & 41,80\% \\
\quad\texttt{\$setuphold} & 7\,314 & 7\,312 & 99,97\% \\
\quad\texttt{\$recrem} & 140 & 140 & 100,00\% \\
\quad\texttt{\$width} & 10\,372 & 0 & 0,00\% \\
\bottomrule
\end{tabular}
\end{center}

A inspeção do SDF encontrou 7\,312 entradas \texttt{SETUPHOLD}, 140
entradas \texttt{RECREM} e nenhuma entrada \texttt{WIDTH}. Esses números
coincidem com os checks anotados. Assim, o percentual global de 41,80\%
é calculado sobre os checks existentes nos modelos Verilog e não indica
que 58,20\% das restrições fornecidas pelo SDF tenham sido rejeitadas.

Os dois checks \texttt{\$setuphold} sem anotação pertencem à instância
\path{testbench.dut.reset_inst_reset_sync_reg_reg[0]}.
A netlist conecta sua entrada D à constante \texttt{1'b1}, o que é
compatível com a ausência de checks de transição de dados no SDF dessa
instância. Os checks de recovery/removal fornecidos foram anotados.

\begin{center}
\begin{tabularx}{\linewidth}{@{}lrX@{}}
\toprule
Aviso & Ocorrências & Significado registrado \\
\midrule
\texttt{NTCNNC} & 132 & Não convergência de valores de checks temporais negativos em determinadas instâncias. \\
\texttt{NTCWID} & 196 & Desativação de checks de largura de pulso porque a diferença entre atrasos NTC de subida e descida supera o limite do check. \\
\texttt{PRLDYN} & 1 & Objetos dinâmicos excluídos da captura de formas de onda pela opção de probing utilizada. \\
\texttt{PRHOPT} & 1 & Sinais otimizados excluídos da captura de formas de onda. \\
\bottomrule
\end{tabularx}
\end{center}

Embora a tabela do anotador informe zero checks na categoria
\emph{Disabled}, o simulador também emite avisos \texttt{NTCWID} de
desativação. As duas informações correspondem a mecanismos distintos e
devem ser preservadas na avaliação. A ausência de mensagens
\texttt{Timing violation} não comprova o atendimento dos checks não
anotados ou desativados, nem resolve os avisos \texttt{NTCNNC}.

\section{Contexto temporal da síntese utilizada}

Como informação complementar, o \path{qor_tsmc28.rpt} da síntese que
gerou a netlist registra período de 10\,000 ps, slack crítico de
$+5\,469{,}7$ ps ($+5{,}4697$ ns), TNS igual a zero e nenhum caminho
violado no grupo \path{pclk_i}. O primeiro caminho de
\path{timing_max_tsmc28.rpt} vai de \path{paddr_i[6]} a \path{pslverr_o}.
A netlist possui 2\,562 células: 706 sequenciais e 1\,856 combinacionais.

Esses valores são resultados da análise estática de síntese, separados
dos resultados de simulação. O fluxo não inclui implementação física,
árvore de clock roteada ou extração de parasitas. Portanto, esta execução
não estabelece fechamento temporal pós-layout nem validação em todos
os corners de processo, tensão e temperatura.

\section{Conclusão}

A regressão pós-síntese TSMC28 a 100 MHz foi aprovada para os estímulos
executados: 398 verificações passaram e as dez operações RSA produziram
os resultados esperados. A configuração sintetizada inclui
\texttt{LOAD\_GAP=1}, e a execução concluiu as operações e a limpeza
de estado e interrupção.

Os atrasos de caminho e todos os checks presentes no SDF foram anotados.
Para ampliar a conclusão temporal, ainda é necessário avaliar os avisos
\texttt{NTCNNC}/\texttt{NTCWID} e a ausência de restrições de largura
de pulso, além das etapas de validação física e de outros corners.

\appendix
\section{Rastreabilidade dos resultados}

Os identificadores completos da Seção 1 definem as variáveis
\texttt{SIM\_RUN} e \texttt{SYNTH\_RUN} usadas nos caminhos abaixo.
São abreviações documentais; substitua-as pelo identificador correspondente.

\begin{itemize}
\item Resumo e entradas da simulação, estatísticas e checks não anotados:
\path{reports/amba_3_apb_rsa_apb_top_tsmc28/xrun_ps/<SIM_RUN>/}, arquivos
\path{simulation_tsmc28.rpt}, \path{simulation_inputs_tsmc28.txt},
\path{sdf_annotation_tsmc28.log} e \path{sdf_unannotated_tsmc28.rpt}.
\item Log completo, com scoreboard e avisos:
\path{logs/amba_3_apb_rsa_apb_top_tsmc28/xrun_ps/<SIM_RUN>/<SIM_RUN>.log}.
\item Netlist e SDF correspondentes:
\path{ps_files/rsa_apb_top_tsmc28/netlist/synth_<SYNTH_RUN>.v} e
\path{ps_files/rsa_apb_top_tsmc28/sdf/synth_<SYNTH_RUN>.sdf}.
\item QoR, área e timing da síntese:
\path{reports/amba_3_apb_rsa_apb_top_tsmc28/genus/<SYNTH_RUN>/}, arquivos
\path{qor_tsmc28.rpt}, \path{area_tsmc28.rpt} e \path{timing_max_tsmc28.rpt}.
\item Configuração elaborada:
\path{logs/amba_3_apb_rsa_apb_top_tsmc28/genus/<SYNTH_RUN>/<SYNTH_RUN>.log}.
\item Formas de onda geradas:
\path{waveforms/amba_3_apb_rsa_apb_top_tsmc28/xrun_ps/<SIM_RUN>/waves_tsmc28.shm}.
\end{itemize}

As formas de onda são indicadas para inspeção posterior; as conclusões
deste relatório foram extraídas dos logs e relatórios textuais.

\end{document}
